\chapter{Sistema de transmisión}
\label{cap:mm-transmision}

\opening{Dos correas, un embrague de rueda libre, tres placas flexibles y dos cajas: todo lo que separa al motor de los rotores.}

\section{Descripción}
\fuente{MM, capítulo 7}

Una polea de aluminio de ranuras múltiples con recubrimiento de acero (la \textbf{polea inferior}) está apernada directamente al cigüeñal. Dos pares de correas en V unidas en banda transmiten la potencia a la \textbf{polea superior}, de aluminio anodizado (opcionalmente con recubrimiento de acero). La polea superior lleva en su cubo un \textbf{embrague de rueda libre tipo sprag} por el que pasa el eje del embrague, que recibe la potencia del sprag y la transmite hacia adelante a la caja del rotor principal y hacia atrás al eje del rotor de cola.

Hay acoplamientos flexibles (placas flexibles) en la entrada de la caja principal y en cada extremo del eje largo del rotor de cola. La caja principal tiene un juego de engranajes cónicos espirales de una etapa, lubricado por salpicadura. El eje de cola no tiene rodamientos colgantes, pero sí un rodamiento amortiguador poco cargado. La caja de cola también tiene engranajes cónicos espirales lubricados por salpicadura. Los ejes de entrada y salida de la caja de cola son de acero inoxidable para resistir la corrosión; los demás ejes son de acero aleado.

Componentes principales: caja y mástil A006; horquilla A006-5 o A908-4 en la entrada de la caja; placa flexible delantera A947-1; horquilla delantera del embrague A907-4 o -5; correas A190; polea inferior A493 y rodamiento inferior A181-4; placa flexible intermedia A947-2; horquilla B223 del eje de cola; eje del rotor de cola B224; placa flexible trasera A947-3; caja de cola A021-1 o B021-1.

\begin{advertencia}
Las placas flexibles A193, sin arandelas pegadas, son obsoletas y deben reemplazarse por placas A947 con arandelas pegadas. Si una arandela pegada se separa de una placa A947, la placa no es aeronavegable y no puede repararse. La placa delantera A947-1 debe ser rev.~E o posterior (letra ``E'' o posterior marcada en dos brazos adyacentes).
\end{advertencia}

\section{Caja de transmisión del rotor principal}

\subsection{Retiro e instalación}

Para retirarla se retiran el rotor principal, el plato oscilante, el carenado del mástil y el capó trasero, el cable del freno rotor, los tanques de combustible (marcando cables), el eje intermedio y sus montantes, y el embrague; se desconectan el detector de virutas, los emisores de efecto Hall, el interruptor del freno y el sensor de sobretemperatura (marcando cables) y la masa; se retiran la manguera y la tobera de enfriamiento y los cuatro pernos de montaje, y se levanta la caja con cuidado para librar los tubos delanteros.

\begin{advertencia}
Puede haber suplementos de nivelación entre los soportes de la caja y las almohadillas del bastidor. Deben quedar en el bastidor o marcarse las almohadillas con su espesor para reinstalarlos.
\end{advertencia}

En la instalación: si los suplementos originales son desconocidos o se instaló un bastidor superior nuevo, nivele las almohadillas; instale pernos y herrajes A650 con su torque; reinstale embrague, eje intermedio, manguera y tobera, tanques y cable del freno (verifique que el interruptor del freno actúe según la figura 7-16A); alinee la costilla inferior del carenado con el eje longitudinal; instale plato oscilante y rotor.

\begin{precaucion}
Antes de operar una caja nueva, con overhaul o reparada, o si alguna vez estuvo acostada horizontalmente, deben inyectarse \textbf{6~onzas de aceite A257-2} por el respiradero de la parte superior trasera del tubo del mástil, manteniendo la caja vertical durante y después de la inyección. Luego llene por el tapón de llenado hasta el nivel indicado en la calcomanía del visor.
\end{precaucion}

\subsection{Nivelación de la caja}

Nivele la aeronave (\S\,18-13) y ponga suplementos A796 en las almohadillas del bastidor de modo que un nivel apoyado lateral o longitudinalmente entre dos almohadillas cualesquiera indique \textbf{$\pm$0,20° o menos}; si no se logra, nivele dentro de un suplemento de 0,020~in. \textbf{Máximo seis suplementos por almohadilla.} Al menos una almohadilla debe quedar sin suplementos (quite igual número de cada una si es necesario).

\subsection{Horquilla de entrada y emisores de efecto Hall}

Para retirar la horquilla A006-5 o A908-4, retire la placa flexible A947-1 y el embrague (o sostenga el eje), ponga un taco de madera entre el cortafuego horizontal y la brida de la horquilla para reaccionar el torque, y retire la chaveta y la tuerca. Para instalarla: arandela A141-10 y tuerca almenada AN320-8 con torque especial y chaveta.

\textbf{Separación emisor de efecto Hall--imán: 0,020 $\pm$ 0,005~in}, ajustable aflojando las contratuercas del emisor; verifique en cada emisor con los pernos-imán, apriete y marque. Verifíquela antes de la prueba de motor o de girar las palas.

\subsection{Sello del piñón}

Con el actuador completamente desacoplado, desconecte la placa intermedia del eje de cola y retire la delantera (anotando los suplementos); apoye el eje del embrague en el cortafuego; retire la horquilla; corte el alambre y retire los pernos de la tapa del rodamiento del piñón y los soportes de los emisores Hall y de sobretemperatura; deslice la tapa fuera del eje.

\begin{precaucion}
No retire los suplementos bajo la tapa del rodamiento del piñón: controlan la precarga del rodamiento y el juego de los engranajes.
\end{precaucion}

Saque el sello con prensa e instale uno nuevo hasta asentarlo, \textbf{con la cara abierta del labio hacia la caja}. Reinstale tapa y soportes, torque especial y alambre de frenado; ajuste la separación de los emisores Hall; reinstale las placas flexibles en su posición original.

\subsection{Anillo O del cárter}

\begin{pasos}
\item Drene la caja y retírela.
\item Mida y registre el juego (\emph{backlash}) del engranaje en un diente, tangencialmente, accesible por el visor o el tapón de llenado.
\item Anote cuidadosamente la posición de cada sujetador, arandela y pila de suplementos en todos los puntos de unión cárter A263-1--carcasa A264-1 (hay igual número de suplementos en cada punto); retírelos e identifíquelos.
\item Retire el cárter y el anillo O. Prelubrique el anillo nuevo A214-172 con el aceite correcto e instálelo sin torcerlo.
\item Arme con los mismos suplementos en las mismas posiciones (los cables de masa van bajo la tuerca detrás del soporte delantero derecho).
\item Torques: ocho tuercas en pernos AN4, \textbf{90~in-lb} (incluido el autobloqueo); dos tornillos MS20074, \textbf{60~in-lb} y alambre; alojamiento del detector de virutas A7260, \textbf{150~in-lb} y alambre; detector A7260, \textbf{75~in-lb} y alambre.
\item Mida de nuevo el juego exactamente igual: debe quedar \textbf{dentro de 0,001~in} del valor registrado.
\item Instale la caja y llénela con el aceite correcto.
\end{pasos}

\subsection{Sobretemperatura de la caja principal}

\begin{itemize}
\item \textbf{Luz \textsf{MR TEMP} encendida y Telatemp con temperatura anormalmente alta:} (1) revise el ducto de enfriamiento (obstrucciones, estado); (2) retire el detector de virutas: si hay virutas, envíe la caja a RHC; (3) drene el aceite, retire visor y tapón y observe los dientes mientras gira el piñón: si hay daño o no gira suavemente, envíe la caja a RHC; si no, rellene; (4) reemplace el Telatemp (sin restos del adhesivo viejo y con buen contacto); (5) si continúan las indicaciones, envíe la caja a RHC.
\item \textbf{Luz encendida pero Telatemp normal:} reemplace el sensor de sobretemperatura y haga los pasos 1, 2 y 5.
\item \textbf{Telatemp en 240~°F/116~°C sin luz:} pruebe el circuito \textsf{MR TEMP} y haga los pasos 1, 2, 4 y 5. La luz se enciende a \textbf{240 $\pm$ 5~°F}.
\end{itemize}

\subsection{Luz de virutas de la caja principal}

Si se enciende \textsf{MR CHIP}: drene y lave la caja colando el aceite (un filtro de pintura sirve) y examine las partículas del aceite y del detector. \textbf{Las partículas de más de 0,12~in de largo o 0,02~in de ancho son motivo de preocupación}: identifíquelas con un imán. Si son ferrosas, envíe la caja a RHC con las partículas; si no, drene y lave la caja. Si la luz vuelve a encenderse dentro de las 100 horas siguientes, puede ser inminente una falla: envíe la caja a RHC.

\section{Conjunto del embrague}

\subsection{Retiro}

\begin{pasos}
\item Retire la placa flexible intermedia.
\item Con la batería conectada, desacople el embrague y ajuste el tornillo de tope inferior del actuador para la máxima holgura de correas; apague la batería. (El actuador debe estar completamente desacoplado antes de retirar el embrague.)
\item Desconecte la placa flexible delantera, el rodamiento superior A184 del actuador, el montante de centrado lateral y el cableado del actuador.
\item Pivote el actuador a la derecha y baje el embrague. Si reinstalará las correas, \textbf{marque su posición delantera/trasera y el sentido de giro}.
\item Retire las correas y el embrague. Si lo envía a RHC, retire y conserve ambas horquillas, el montante y la ferretería.
\end{pasos}

\subsection{Rodamiento superior A184}

Se retira con prensa y el extractor inferior MT303-5, o en campo con el extractor superior MT303-2, el tapón MT303-16 y el pasador MT303-13 (roscas con antiagarrotante); cuide que el embrague no caiga.

\begin{advertencia}
El rodamiento A647-4 del conjunto A184-3 tiene una holgura interna modificada por RHC. NO LO SUSTITUYA.
\end{advertencia}

Instalación: limpie el eje; caliente el rodamiento (por ejemplo, con una lámpara a un pie) \textbf{sin superar 180~°F} según su Telatemp; aplique imprimante de cromato de zinc o epóxico, o adhesivo B270-10, en el asiento y, húmedo, instale el rodamiento con la oreja del montante de centrado a la izquierda y la unión del actuador abajo, presionando sobre la pista interior con el tubo MT303-19 (o la prensa de campo MT303-4) hasta la cota de la figura 7-3D (1,480 a 1,490~in). \textbf{No apoye el embrague en la polea al presionar ni presione más allá de la cota.} Verifique que gire suavemente, selle con imprimante la unión pista--eje, marque la pista con franjas en dos puntos a 180° y reemplace el Telatemp si se alteró con el calor.

\subsection{Sellos del embrague}

\begin{pasos}
\item Retire el embrague (y el rodamiento A184-3 si cambia el sello trasero). Oriente el embrague vertical con el sello que pierde arriba.
\item Retire los seis pernos del círculo interior del retenedor. Si el retenedor es circular con dos escotaduras de media luna, retire los cuatro pernos exteriores pero \textbf{no los dos con arandelas de gran área}, que mantienen armado el embrague interior (no se aprueba más desarme). Si tiene dos tornillos de cabeza hueca, retire los seis exteriores y esos dos tornillos.
\item Retire el retenedor con el sello y descarte el anillo O de la carcasa. Saque el sello.
\item \textbf{Si el alojamiento del sello mide más de 1,687~in, reemplace el retenedor.}
\item Sellos ``National'': cubra el alojamiento con sellador B270-3 y presione a temperatura ambiente. Sellos ``CR'' (exterior verde): caliente el retenedor a 150~°F y presione en seco. El sello queda al ras del labio exterior.
\item Llene la cavidad del sprag y los rodamientos con aceite A257-4 hasta el ras de las pistas, girando el eje.
\item Lubrique levemente con A257-4 el anillo O nuevo y el labio del sello; instale el retenedor (las escotaduras sobre las dos arandelas A141-14), con las cabezas de los pernos hacia el extremo corto del eje; apriete en cruz con torque estándar y marque.
\end{pasos}

\subsection{Instalación del embrague}

Instale las horquillas A195 y A907 si se retiraron, conecte el montante de centrado al rodamiento A184, verifique el actuador desacoplado e identifique el tipo de poleas (ranura profunda o baja). Monte el embrague, pase las correas sobre la horquilla trasera \textbf{en su orden y sentido correctos}, conéctelo al actuador, instale la placa A947-1 y el montante de centrado al bastidor. Verifique la alineación de poleas, instale la placa intermedia, ajuste el tope inferior del actuador para la deflexión de correas especificada, haga la prueba de motor (\S\,2.210), reajuste y apriete la contratuerca del tope, y marque las franjas.

\begin{precaucion}
RHC no recomienda usar poleas superiores A170-1 o -2 rev.~N o anteriores con poleas inferiores A493-1 rev.~T o posteriores: pueden causar vibración excesiva de las correas. Profundidad de ranura de la polea superior: 0,92~in (``ranura profunda'', rev.~N y anteriores) o 0,71~in (``ranura baja'', rev.~O y posteriores).
\end{precaucion}

\section{Alineación de las poleas}

\begin{pasos}
\item Verifique que el suplementado del motor esté dentro de límites (capítulo~\ref{cap:mm-motor}).
\item Acople el embrague (ventilador y carcasa instalados) y verifique el ángulo del eje del embrague.
\item Apoye la barra de alineación MT331-4 contra la cara trasera de la polea inferior, extendida hacia arriba hasta la línea central horizontal de la polea superior, y mida las separaciones izquierda y derecha.
\end{pasos}

\begin{table}[htbp]
\centering\small
\caption{Separaciones de alineación de poleas (barra MT331-4 a la polea superior).}
\begin{tabularx}{\linewidth}{@{}Ycc@{}}
\toprule
Embrague & Separación izquierda & Separación derecha \\
\midrule
A018-1, polea superior anodizada (caras lisas) & 0,030 a 0,070 in & 0,010 a 0,050 in \\
A018-2, polea superior metalizada (caras rugosas) & 0,010 a 0,050 in & 0,001 a 0,030 in \\
\midrule
\multicolumn{3}{@{}p{\linewidth}@{}}{En ambos casos, la izquierda debe ser 0,010 a 0,030~in mayor que la derecha.} \\
\bottomrule
\end{tabularx}
\end{table}

Si la izquierda no es 0,010--0,030~in mayor, ajuste el centrado lateral o los suplementos del motor. \textbf{Acortar el montante de centrado aumenta la separación izquierda y disminuye la derecha}; alargarlo, lo contrario.

\textbf{Si alguna separación excede el máximo}, el embrague está muy adelante o el motor muy atrás: adelante el motor quitando suplementos iguales de cada soporte (respetando el mínimo), o atrase el embrague suplementando la placa flexible delantera (máximo \emph{una} arandela AN960-516L entre la placa A947-1 y ambos brazos de la horquilla A907 y/o A908; cada arandela reduce ambas separaciones en su espesor) o alargando la horquilla A907. Las A907-4 y -5 tienen dos juegos de agujeros que cambian su largo efectivo en 0,120~in; la A907-5 es 0,120 o 0,240~in más larga que la posición larga de la A907-4. Cada 0,120~in de alargue reduce ambas separaciones 0,120~in.

\textbf{Si alguna separación es menor que el mínimo}, el embrague está muy atrás o el motor muy adelante: atrase el motor agregando suplementos iguales (respetando el máximo), o adelante el embrague quitando arandelas de la placa delantera o acortando la horquilla A907.

\begin{precaucion}
Desacople el embrague antes de suplementar la placa A947-1. La horquilla A907-4 reemplaza a las A907-1 y -2; la A907-5 reemplaza a la A907-3 e incluye un juego de agujeros adicional.
\end{precaucion}

Luego: verifique el suplementado de la placa intermedia si se cambiaron suplementos en la delantera o la posición de la A907; ajuste la guía de correas (separación de \textbf{0,18~in con el embrague acoplado}, al ras de la cara delantera de la polea, y 0,18~in entre el tubo de empuje--tracción y los pernos de la guía; en los bastidores A020-2 S/N 0501 y posteriores el soporte de la guía va soldado), y verifique la correlación del acelerador.

\subsection{Ángulo del eje del embrague}

Con las palas perpendiculares al eje del helicóptero, ponga un transportador sobre el cubo del rotor principal mirando a la derecha y llévelo a cero; acople el embrague y mida el ángulo del eje (con el transportador mirando en el mismo sentido). \textbf{Límites con el embrague acoplado: 2,5° a 5,0°.} La placa flexible no tiene desalineación vertical a 3,5°. Si está fuera de límites, deben reemplazarse las correas.

\subsection{Centrado lateral del embrague}

Con el actuador completamente acoplado, verifique las separaciones (la izquierda 0,010--0,030~in mayor). El montante puede ajustarse con el embrague acoplado siempre que se mantengan: \textbf{mínimo 0,030~in entre el motor del actuador y el tubo del bastidor superior}, y \textbf{mínimo 0,150~in entre las correas y el protector derecho del bastidor} (actuador acoplado). Si no se logra, consulte a RHC. Torque estándar de pernos y contratuercas.

\subsection{Horquillas A907 y A195 del eje del embrague}

Se retiran con prensa (botador de bronce o aluminio contra el borde exterior del eje, no contra el espaciador interior) o, la A907, girando el eje con la horquilla sujeta suavemente en una morsa acolchada después de aplicar aceite penetrante; la A195 con los extractores MT303. Los brazos de la A195 se doblan con facilidad: retírela siempre para transportar el embrague. Para instalarlas: decape y limpie el asiento del eje, el alojamiento de la horquilla y las zonas de los bloques de apriete; cubra con imprimante aprobado (de preferencia cromato) y, húmedo, instale la horquilla (alineando los agujeros marcados) y los bloques imprimados; torque estándar, \emph{palnuts}, imprimante con filetes en todas las uniones y franjas. \textbf{No use ningún otro lubricante.}

\section{Correas en V}

\subsection{Retiro}

Retire el ventilador; desacople el actuador y ponga un taco de madera de 6~in entre el eje del embrague y el cortafuego horizontal (evita la caída de la polea y protege la placa delantera); desconecte el montante de centrado y retire la placa intermedia; marque las correas (delantera/trasera y sentido de giro); afloje y corra la guía; desconecte el actuador y saque las correas sobre la horquilla trasera.

\begin{precaucion}
Las correas usadas deben reinstalarse en su orden original, por sus diferencias individuales de estiramiento.
\end{precaucion}

\subsection{Instalación}

\begin{pasos}
\item Evalúe la correa del alternador e inspeccione (y lubrique) los rodamientos del actuador.
\item Identifique el tipo de poleas e inspecciónelas: reemplace las que tengan corrosión, picaduras o descascarado del recubrimiento, desgaste del anodizado, surcos, asperezas o crestas filosas (pueden hacer que las correas rueden, se rompan o se salgan).
\item Pinte las ranuras de la polea inferior con una capa delgada de imprimante aprobado (de preferencia cromato; puede hacerse sin retirarla, con aerosol o aerógrafo).
\item Inspeccione las correas (capítulo~\ref{cap:mm-inspecciones}) e instálelas en su orden y sentido. \textbf{Las correas son un juego apareado: instálelas solo como juego.}
\item Conecte el actuador y el montante (torque estándar) e instale ventilador y carcasa.
\item \textbf{Deflexión con el actuador desacoplado: 1,4 a 1,6~in hacia adentro con presión del dedo} (1,25~in mínimo con correas nuevas; 1,4--1,6~in nominal después del asentamiento). Ajuste el tornillo de tope inferior y apriete su contratuerca.
\item Verifique la alineación de poleas y suplemente la placa intermedia.
\item Gire la transmisión a mano con el actuador desacoplado y verifique una holgura mínima de \textbf{0,125~in con correas nuevas (0,25~in después de unas 10~h)} entre el acoplamiento intermedio y el perno de unión balancín A331-1--tubo A121-17. Si es menor, suplemente el motor.
\item Posicione la guía de correas y, si es necesario, instale la abrazadera protectora del bastidor superior.
\item Haga funcionar en tierra o en estacionario al menos media hora y revise el patrón de desgaste del imprimante en las cuatro ranuras inferiores: debe ser similar en todas (vea el capítulo~\ref{cap:mm-inspecciones}). Si no lo es, verifique la alineación, reemplace las correas y repita; si persiste, puede ser necesario reemplazar la polea inferior. \textbf{El patrón debe ser similar en las cuatro ranuras antes de liberar la aeronave al vuelo.}
\item Verifique el ángulo del eje del embrague.
\end{pasos}

\begin{precaucion}
Las correas demasiado tensas pueden dañar las placas flexibles en el arranque; las demasiado sueltas pueden saltar de las ranuras al acoplar. Una demora de más de 5~s entre el acople y el giro del rotor indica holgura excesiva; si el rotor gira al dar arranque al motor, puede faltar holgura. Con correas nuevas, el tope inferior debe ajustarse con más frecuencia.
\end{precaucion}

\section{Eje del ventilador, soporte de la corona, polea inferior y correa del alternador}

Retire el ventilador y la carcasa (pueden salir sin partir la carcasa); apoye el eje del embrague con un taco; retire el perno inferior actuador--rodamiento A181-4; marque las correas; corte el alambre y retire pernos y arandelas; retire el eje del ventilador (y la polea inferior golpeándola con un mazo blando mientras tira). Para cambiar la correa del alternador o acceder a la nariz del motor hay que retirar el soporte de la corona: afloje la correa del alternador y marque el buje del cigüeñal que coincide con la marca ``O'' del soporte.

En la instalación, instale la correa nueva y la de repuesto si corresponde y el soporte alineando la marca ``O'' con el buje marcado (ese buje y su agujero son levemente mayores). Instale la polea inferior y el eje con seis pernos NAS6606H23 y arandelas NAS1149F0632P (o NAS6606H24 con NAS1149F0632P y NAS1149F0663P), con torque especial en cruz y frenados por pares con alambre inoxidable de 0,041~in. Reponga las correas en su orden, el perno inferior del actuador, el ventilador, retire el taco y ajuste la tensión del alternador (Lycoming SI 1129).

\begin{advertencia}
Los rodamientos de los ejes de ventilador A007-3 y A007-5 tienen holguras internas modificadas por RHC. NO SUSTITUYA ningún otro rodamiento.
\end{advertencia}

\section{Eje del rotor de cola}

\subsection{Retiro e instalación}

Desconecte el eje de la placa intermedia (dejándola en la horquilla del embrague) y apóyelo con un taco de 4~in cerca de su extremo delantero (protege la placa trasera); retire el cono de cola; desconecte los dos pernos NAS1303 del soporte del eje (con una extensión de 3~ft y dado de 3/8~in) y el tubo de empuje--tracción del balancín de la caja (si no se cambia su largo, no hay que reglar de nuevo); retire los cuatro pernos de la caja de cola (un NAS1304-8 y tres MS20074-04-06) y retírela lo suficiente para soltar los pernos de la placa trasera; marque eje y placa y desconéctelos; saque el eje por el frente del cono.

En la instalación, \textbf{las arandelas pegadas de la placa flexible trasera deben quedar directamente contra la brida del eje de entrada de la caja} (una placa mal instalada puede dañar el eje y la caja). Los MS20074-04-06 se aprietan con su torque y se frenan con alambre de 0,032~in. Conecte el amortiguador al bastidor del cono y, por los agujeros de inspección, verifique la alineación longitudinal de su grapa (o brazo) con el soporte del cono: \textbf{si la holgura o interferencia supera 0,12~in, consulte a RHC}; verifique su orientación. Instale el cono, el tubo A121-17, la placa intermedia con su suplementado y verifique el descentramiento.

\subsection{Fricción del amortiguador}

Desarme el amortiguador e inspeccione las arandelas DU A041-5 (recubrimiento de teflón gastado: cara gris oscuro), las arandelas A141-15 y el brazo A041-3 (desgaste o surcos). Rearme con las arandelas elásticas con la cara cóncava hacia las arandelas DU, y \textbf{la cara de teflón (gris oscuro) de la DU contra la arandela A141-36 o el brazo A041-3} (retire cualquier recubrimiento plástico o cinta). Mida la fricción de pivote con una balanza de resorte o peso muerto: \textbf{1,5 a 2,5~lb} para mover la grapa medida en el punto A (girando en torno a B) y \textbf{1,5 a 2,0~lb} medida en B (girando en torno a C). Si es baja, doble la arandela elástica A041-6 hasta \textbf{0,070 $\pm$ 0,010~in} de altura total; si aún es baja, rebaje el extremo del muñón A105-9. Si es alta, aplane levemente la arandela o reemplace el muñón.

\begin{precaucion}
No ajuste la fricción del amortiguador cambiando el torque del perno.
\end{precaucion}

\subsection{Instalación y suplementado de la placa flexible intermedia}

Mediciones con la placa intermedia retirada y el embrague acoplado:

\begin{pasos}
\item Gire las horquillas del eje de cola y A195 hasta que sus brazos queden horizontales.
\item Inserte un perno NAS1304 por ambas horquillas a las 9 en punto y mida la separación entre brazos a las 3; luego inserte el perno a las 3, retire el de las 9 y mida a las 9 (medición 1).
\item Gire el eje 180° y repita (medición 2). Ambas deben ser similares; si no, una o ambas horquillas están dobladas.
\item La diferencia entre las separaciones de las 3 y las 9 \textbf{no debe exceder 0,050~in} en ninguna medición; si la excede, ajuste el montante de centrado y alinee las poleas.
\item Con la medición de menor diferencia: \emph{dimensión calculada} = (lectura de las 9 + lectura de las 3) / 2 $-$ espesor de la placa A947-2 (medido sobre las arandelas pegadas en cuatro puntos y promediado).
\end{pasos}

\begin{table}[htbp]
\centering\footnotesize
\caption{Suplementado de la placa flexible intermedia (MM tabla 7-1).}
\begin{tabularx}{\linewidth}{@{}lYY@{}}
\toprule
Dimensión calculada (in) & Lado embrague: entre A947-2 y A195-1 (2 lugares) & Lado eje de cola: entre A947-2 y B223 (2 lugares) \\
\midrule
0,135 o más & \multicolumn{2}{p{5.2cm}}{Quitar suplementos de la caja de cola, suplementar la placa A947-1\textsuperscript{**}, cambiar el largo de la A907\textsuperscript{**} o instalar una B223 más larga} \\
0,134 a 0,105 & AN960-416 & AN960-416 \\
0,104 a 0,075 & AN960-416 & AN960-416L \\
0,074 a 0,045 & AN960-416L & AN960-416L \\
0,044 a 0,015 & AN960-416L & Ninguna \\
0,014 a $-0,015$ & Ninguna & Ninguna \\
$-0,016$ o más negativo & \multicolumn{2}{p{5.2cm}}{Suplementar la caja de cola hacia atrás (máx. 0,080~in\textsuperscript{*}), quitar suplementos de la A947-1\textsuperscript{**}, cambiar el largo de la A907\textsuperscript{**} o instalar una B223 más corta} \\
\bottomrule
\multicolumn{3}{@{}p{\linewidth}@{}}{\textsuperscript{*} Al suplementar la caja de cola, reemplace los pernos MS20074-04-04 por MS20074-04-06; con más de 0,040~in de suplementos, cambie las arandelas bajo las cabezas a AN960-416L. El total no debe exceder 0,080~in. \textsuperscript{**} Requiere verificar la alineación de poleas.}
\end{tabularx}
\end{table}

Al agregar o quitar suplementos entre la caja de cola y la fundición del cono, \textbf{alargue el tubo A121-17 media vuelta del terminal por cada 0,018~in} que la caja se desplace hacia atrás (y acórtelo igual al quitarlos). Instale la placa con los suplementos determinados, torque, \emph{palnut} y franja (arandelas bajo las tuercas para dejar 2 a 4 hilos). Verifique con el actuador desacoplado una holgura mínima de \textbf{0,250~in} entre la placa y el perno del terminal delantero del tubo A121-17 en el balancín A331.

\subsection{Descentramiento del eje de cola}

Un descentramiento excesivo puede hacer fallar el acoplamiento intermedio o el amortiguador. Retire las tapas de inspección del lado derecho del cono, acople el embrague, arme la herramienta MT260-6 con un comparador (sirve el del medidor de estiramiento de pernos MT122), introdúzcalo por el agujero de inspección más trasero apoyando firmemente la carcasa contra el cono y con la extensión sobre el eje, y haga girar la transmisión (desde el eje del embrague A166) al menos tres vueltas completas, promediando. Repita en cada agujero. \textbf{El descentramiento máximo en cualquier punto es 0,025~in}; si se excede, el eje debe repararse o reemplazarse.

\subsection{Eje de cola de dos piezas}

Consta del eje B224-1, la horquilla B223-1 a -4 (o -6), el rodamiento amortiguador A041-1, el bloque A553-2 y dos bloques A207-4. El amortiguador solo es reemplazable en campo con este eje. El B224-1 se fabrica en un solo largo; el largo total se ajusta con horquillas B223 en incrementos de 0,150~in: -6, 0,750~in; -1, 0,900; -2, 1,050; -3, 1,200; -4, 1,350. Equivalencias con los ejes de una pieza: A197-7 (140,25~in) usa B223-1; A197-4 (140,40) B223-2; A197-5 (140,55) B223-3; A197-6 (140,70) B223-4; A197-1 según se requiera.

Si el largo es desconocido: instale la placa intermedia en la horquilla A195-1 sin suplementos; instale el bloque A553-2 en el extremo delantero del eje; conecte el eje a la placa trasera y la caja, instale la caja, el amortiguador y el cono, y acople el actuador. Con el eje centrado en la placa, mida desde la placa hasta el centro del agujero delantero del eje y reste 0,032~in (espesor de la arandela): esa es la dimensión calculada. Elija la B223 más cercana sin excederla; la diferencia se compensa con suplementos. Instale la horquilla con dos pernos NAS1304-21 con arandela y bloque A207-4 bajo cabeza y tuerca, con \textbf{100~in-lb más arrastre (110~in-lb en total)} y \emph{palnuts}; selle las uniones con imprimante. El rodamiento se instala en el eje con adhesivo B270-10 (brida del alojamiento hacia el extremo soldado del eje), a \textbf{96,59 $\pm$ 0,03~in}, sellando la unión con imprimante.

\subsection{Eliminación de corrosión del eje B224-1}

\begin{itemize}
\item No se permiten herramientas eléctricas; retire solo el material mínimo.
\item \textbf{Máximo 0,001~in de remoción en cualquier punto de la superficie (0,002~in de reducción del diámetro)}; espesor mínimo de pared de \textbf{0,032~in} en la zona trabajada.
\item Lije o pula paralelo al eje longitudinal; termine con lija al agua de grano 400 o más fino e imprima el metal desnudo con cromato de zinc o epóxico.
\end{itemize}

\section{Caja del rotor de cola}

\subsection{Retiro e instalación}

Marque cubo, palas y varillas de paso respecto de sus brazos de control antes de desarmar (así no hay que reglar si no se cambian los largos de los tubos). Retire el rotor de cola dejando las varillas en las palas; desconecte el balancín A120-1 del tubo y de la caja y saque el control de paso con el balancín por el eje de salida; desconecte el detector de virutas; retire el eje de cola y la placa flexible trasera. Para instalar: eje de cola; control de paso y balancín (verifique su suplementado, \S\,8.562); tubo A121-17; rotor de cola.

\subsection{Suplementado de la caja de cola}

Solo se requiere cuando la separación entre las horquillas del embrague y del eje de cola es demasiado pequeña para instalar la placa intermedia. \textbf{Máximo 0,080~in hacia atrás.} Con suplementos use tres pernos MS20074-04-06 y uno NAS1304-8; con más de 0,040~in, cambie las arandelas NAS1149F0463P bajo las cabezas por NAS1149F0432P. Con el sistema ya reglado, extienda el terminal trasero del A121-17 media vuelta por cada 0,018~in de suplemento (y retráigalo al quitarlos).

\subsection{Luz de virutas de la caja de cola}

Si se enciende \textsf{TR CHIP}: drene y lave la caja colando el aceite y examine las partículas. \textbf{Las de más de 0,09~in de largo o 0,02~in de ancho son motivo de preocupación}: si son ferrosas, envíe la caja a RHC con ellas; si no, drene y lave. Si la luz vuelve a encenderse dentro de las 100 horas siguientes, envíe la caja a RHC.

\subsection{Sellos de la caja de cola}

\textbf{Sello de salida de la A021:} drene, limpie el detector y el visor (reinstálelos con torque especial y alambre); retire el rotor y el control de paso; verifique el torque de los pernos cono--bastidor; deslice el mandril MT295-10 hasta el sello y fíjelo con cinta; perfore cuatro agujeros en el sello con el punzón MT295-3 guiado por el mandril (no use otras herramientas: podría dañar la carcasa, el eje o los rodamientos); gire el mandril, inserte tornillos 92470A113 (\textbf{no más largos de 0,75~in}) en el sello y extráigalo con el tornillo MT295-6 lubricado. Limpie los asientos (sin pulirlos ni alterarlos; si hay daño, envíe la caja a RHC), lubrique con el aceite correcto el labio y el asiento del eje (no el borde exterior, que va limpio y seco) y deslice el sello nuevo A966-2 con la cara cóncava hacia la caja; instálelo al ras (tipo enrasado) o hasta que el mandril toque la caja (tipo rebajado). Llene, verifique los respiraderos del tapón, apriete y reinstale control y rotor; revise fugas después de balancear.

\textbf{Sello de salida de la B021:} retire los cuatro pernos de cabeza perforada del retenedor B552-1 y sáquelo \textbf{sin alterar los suplementos entre retenedor y rodamiento} (precarga). Presione el sello nuevo A966-4 con el labio hacia la caja \textbf{0,140 a 0,160~in por debajo} de la superficie plana exterior; anillo O nuevo lubricado; lubrique el asiento del eje; instale el retenedor, los pernos con torque especial y alambre.

\textbf{Sello de entrada de la B021:} drene; retire el cono de cola, el eje de cola, la placa trasera y la horquilla estriada B546 (tuerca y arandela); corte el alambre y retire el retenedor B549 y el suplemento A117; descarte sello y anillos O; retire el espaciador A266. Inspeccione: espaciador sin surcos, suplemento sin rayado y estrías sin escalones de desgaste. Cubra las estrías con protector B270-21; presione el sello A966 nuevo \textbf{0,030~in por dentro del ras} en el extremo interior achaflanado del retenedor (lado cerrado hacia el extremo hexagonal); anillos O nuevos prelubricados en el retenedor y en el eje; suplemento A117 hasta la pista exterior; espaciador A266 lubricado con el chaflán interior primero. Retenedor con antiagarrotante A257-9 en las roscas, torque especial y alambre; horquilla con adhesivo B270-11 en las roscas del eje y, húmedo, tuerca con torque especial y franja.

\section{Actuador del embrague}

\subsection{Retiro e instalación}

Retire el carenado del cono, desacople por completo, retire el ventilador y ponga un taco de 4~in bajo el eje del embrague, delante de la polea superior, sobre el cortafuego horizontal; desconecte los dos conectores y retire los pernos NAS6604 superior e inferior. Para instalarlo: conecte al rodamiento A184 (torque estándar, \emph{palnut}, franja); perno inferior con la cabeza hacia atrás, arandela NAS1149F0432P bajo la cabeza y NAS1149F0463P bajo la tuerca; retire el taco, instale el ventilador y conecte el motor y los interruptores.

\begin{precaucion}
No acople el actuador sin la carcasa instalada.
\end{precaucion}

\subsection{Reemplazo del motorreductor}

Puede hacerse con el actuador instalado. Desconecte los cables, corte el alambre y retire los cuatro tornillos (sin dejar caer nada en las poleas) y saque el motorreductor (puede requerir un movimiento de desenroscado). Lubrique levemente el tornillo sin fin del nuevo con grasa A257-1 y ponga el anillo O en la nariz. Instálelo con los cables hacia abajo, \textbf{sin separación entre la carcasa y la brida antes de poner los tornillos} (no use los tornillos para arrastrarlo: la brida puede romperse); frene con alambre de 0,020~in por el método de hebra simple. Conecte, acople el embrague escuchando trabas y verifique que tense por completo las correas; selle sujetadores y juntas con B270-8 y amarre los cables.

\begin{precaucion}
Si los cables del motorreductor se instalan invertidos en el conector, el motor funcionará al revés y los interruptores de tope inferior y de resortes no cortarán la corriente, con daños graves al actuador y a las correas.
\end{precaucion}

\subsection{Ajustes del actuador}

\begin{itemize}
\item \textbf{Límite acoplado:} lo determinan los resortes de columna, que controlan la tensión de las correas; \textbf{solo la fábrica puede ajustarlos}.
\item \textbf{Tope inferior (desacoplado):} lo activa el tornillo de tope inferior, ajustable con llaves largas de 3/8~in (boca) y 1/4~in (corona), para mantener la deflexión de correas especificada. Separación mínima entre cabezas de tornillos en la tijera: \textbf{0,005~in}.
\item \textbf{Extensión máxima acoplada: 1,85~in}, medida entre la parte superior del bloque de resortes de columna y la parte inferior de la carcasa del tornillo. \textbf{Al alcanzarla, deben reemplazarse las correas.}
\end{itemize}

\subsection{Interruptores y fusible}

Los interruptores pueden reemplazarse como arnés completo o individualmente, soldando y cubriendo con termorretráctil. Llevan dos arandelas espaciadoras entre interruptor y placa en cada tornillo y una bajo la tuerca; no apriete en exceso (la carcasa plástica puede agrietarse). \textbf{El fusible en línea del motor del actuador, montado en el panel de interruptores de prueba, es de 1,5~A (AGC 1,5).}

\section{Freno rotor}

\textbf{Retiro:} retire el capó del cono y la placa flexible delantera (anote los suplementos), la chaveta, la tuerca y la horquilla de entrada de la caja; desconecte el resorte, el cable y el microinterruptor (conector de tres pines); corte el alambre y retire los tres pernos.

\textbf{Instalación:}
\begin{pasos}
\item Monte el freno sobre el eje estriado del piñón con dos espaciadores cortos A130-41 entre freno y tapa del piñón, y los tres pernos NAS1352-4H30P con arandela NAS1149F0432P (uno por el freno y un espaciador; otro por la placa del emisor Hall, el freno y un espaciador; el tercero por la placa del emisor Hall y el espaciador largo). Torque especial y alambre de 0,032~in; enganche el resorte.
\item Cubra las estrías con protector B270-21; instale la horquilla A908-4, arandela A141-10 y tuerca AN320-8 con torque especial hasta alinear la ranura, \textbf{sin exceder 410~in-lb}; chaveta MS24665-210.
\item Ajuste la separación de los emisores Hall; conecte cable y microinterruptor.
\item Con el freno aplicado, ajuste la separación entre la cabeza del perno NAS428-3-12 del brazo B109-1 y el espaciador A130-43 a \textbf{0,030--0,035~in} y trábela con la tuerca.
\item Con el freno suelto, la separación entre la palanca B112-1 y el brazo B110-1 debe ser \textbf{0,030--0,170~in}; ajústela con el largo de la cadena de bolitas (una bolita ≈ 0,150~in).
\item Ajuste el microinterruptor para abrir y cerrar con \textbf{0,20--0,30~in} entre la palanca B112-1 y el brazo B110-1.
\item Instale la placa delantera, verifique la alineación de poleas e instale la placa intermedia.
\end{pasos}

\textbf{Pastillas:} espesor mínimo \textbf{0,030~in}. Se reemplazan reemplazando los conjuntos de brazo B110-1 y B109-1.
